\documentclass[tikz,border=12pt]{standalone}
\usepackage{fontspec}
\usepackage{xeCJK}
\IfFontExistsTF{Segoe UI}{\setmainfont{Segoe UI}}{\setmainfont{Arial}}
\IfFontExistsTF{Microsoft JhengHei}{\setCJKmainfont{Microsoft JhengHei}}{\setCJKmainfont{Noto Sans CJK TC}}
\xeCJKsetup{CJKglue={\hskip 0.09em plus 0.02em}, CJKecglue={\hskip 0.12em plus 0.02em}}
\usetikzlibrary{arrows.meta,calc,positioning}
\definecolor{paper}{HTML}{FFFDF8}
\definecolor{ink}{HTML}{0F172A}
\definecolor{muted}{HTML}{475569}
\definecolor{line}{HTML}{CBD5E1}
\definecolor{blue}{HTML}{2563EB}
\definecolor{bluebg}{HTML}{EAF2FF}
\definecolor{green}{HTML}{00856A}
\definecolor{greenbg}{HTML}{E6FAF2}
\definecolor{gold}{HTML}{D97706}
\definecolor{goldbg}{HTML}{FFF2C7}
\definecolor{red}{HTML}{DC2626}
\definecolor{redbg}{HTML}{FFE8E8}
\tikzset{
  title/.style={font=\fontsize{16.5}{22}\selectfont\bfseries,text=ink,align=center},
  sub/.style={font=\fontsize{9.8}{14}\selectfont,text=muted,align=center},
  chip/.style={rounded corners=12pt,line width=.9pt,align=center,inner xsep=12pt,inner ysep=8pt,font=\fontsize{9.5}{13}\selectfont,text=ink},
  row/.style={rounded corners=11pt,line width=.9pt,align=left,inner xsep=14pt,inner ysep=8pt,font=\fontsize{9.6}{13.2}\selectfont,text=ink},
  label/.style={font=\fontsize{9}{12}\selectfont\bfseries,text=ink,align=center},
  small/.style={font=\fontsize{8.8}{12}\selectfont,text=muted,align=center},
  arrow/.style={-{Stealth[length=3mm,width=1.9mm]},draw=ink,line width=1.0pt}
}
\begin{document}
\begin{tikzpicture}[x=1cm,y=1cm,line cap=round,line join=round]
\path[fill=paper] (0,0) rectangle (16,9);

\node[title,text width=13.8cm] at (8,8.16) {不要把整張桌子掃進 prompt};
\node[sub,text width=12.4cm] at (8,7.5) {長脈絡的問題不是容量，而是整理後能不能被下一輪使用。};

\draw[rounded corners=20pt,draw=red,line width=1.05pt,fill=redbg] (0.9,1.45) rectangle (5.75,6.55);
\node[label] at (3.32,5.95) {雜訊抽屜};
\node[chip,fill=white,draw=line,text width=2.8cm] at (3.32,4.92) {整段對話};
\node[chip,fill=white,draw=line,text width=2.8cm] at (3.32,3.85) {所有錯誤};
\node[chip,fill=white,draw=line,text width=2.8cm] at (3.32,2.78) {過期假設};
\node[small,text width=3.6cm] at (3.32,1.93) {看似完整\\其實很難接手};

\draw[arrow] (6.05,4.05)--(7.25,4.05);
\node[small,text width=1.25cm] at (6.65,4.55) {整理};
\node[small,text width=1.25cm] at (6.65,3.55) {刪減};

\draw[rounded corners=20pt,draw=green,line width=1.05pt,fill=greenbg] (7.55,1.15) rectangle (15.1,6.85);
\node[label] at (11.32,6.18) {脈絡帳本};
\node[row,fill=bluebg,draw=blue,text width=5.65cm] at (11.32,5.18) {\textbf{任務狀態}：做到哪一步，誰接手};
\node[row,fill=goldbg,draw=gold,text width=5.65cm] at (11.32,4.1) {\textbf{可重用規則}：下次直接檢查};
\node[row,fill=white,draw=line,text width=5.65cm] at (11.32,3.02) {\textbf{待查疑點}：不能假裝已確認};
\node[row,fill=redbg,draw=red,text width=5.65cm] at (11.32,1.94) {\textbf{不要帶入}：個資與過期例外};

\node[small,text width=11.6cm] at (8,.58) {context engineering 不是讓模型背更多，而是讓資訊各自有位置。};
\end{tikzpicture}
\end{document}
